\documentclass[12pt, a4paper, notitlepage]{article}
\usepackage[utf8]{inputenc}

%math
\usepackage{amsmath, amssymb, amsthm}

%graphics
\usepackage{graphicx, float}
\graphicspath{{images/}}

%customisation
\usepackage[top=2.5cm, bottom=2.5cm, left=2.5cm, right=2.5cm, heightrounded]{geometry}
\usepackage{ragged2e}
\renewcommand{\baselinestretch}{1.15}
\setlength{\parindent}{0pt}
\setlength{\parskip}{0.8em}
\usepackage{newtxtext}
\usepackage{newtxmath}
\renewcommand{\figurename}{Slika}

\usepackage[backend=biber,style=numeric]{biblatex}
\addbibresource{Viri.bib}


\usepackage{fancyhdr}
\setlength{\headheight}{15pt}
\pagestyle{fancy}
\fancyhf{}

\fancyhead[L]{PROGRAMIRANJE C++}
\fancyhead[C]{Program organizacije prostorov}
\fancyhead[R]{\textit{ \nouppercase{\leftmark}}}
\fancyfoot[C]{\thepage}

\fancypagestyle{plain}{ %
  \fancyhf{} % remove everything
  \renewcommand{\headrulewidth}{0pt} % remove lines as well
  \renewcommand{\footrulewidth}{0pt}
}

\usepackage{datetime}

%%%%%%%%%%%%%%%
%% BEGINNING %%
%%%%%%%%%%%%%%%

\setcounter{page}{0}
\author{Loris Bovcon}
\begin{document}

  \begin{titlepage}

    \flushleft
    Škofijska gimnazija Vipava\\
    Goriška cesta 29\\
    5271 Vipava
    \vspace{1cm}

    \begin{center}
      \large
      Seminarska naloga pri predmetu informatika\\
      \rule{\textwidth}{1pt}
      \vspace{0cm}

      \Huge
      \textbf{
        PROGRAMIRANJE V C++\\
      }
      \Large
      \vspace{0.5cm}
      Program organizacije prostorov
    \end{center}

    \flushright
    \vfill
    Avtor: Loris Bovcon 4.b\\Mentor: Bogdan Urdih, prof.
    \vspace{1cm}
    \center
    Ajdovščina, \the\day.\the\month.\the\year

  \end{titlepage}

  \newpage
  \section{Povzetek}
  
  \newpage
  \tableofcontents
  \newpage

%%%%%%%%%%%%%%%%
%%% CONTENTS %%%
%%%%%%%%%%%%%%%%

  \justifying
  \section{Uvod}
  \subsection{Cilji}
  Cilji seminarske naloge so bili sledeči:
  \begin{enumerate}
    \item narediti program, ki bi mi pomagal organizirati sobo in najti stvari, če pozabim kje so,
    \item utrditi predhodnje znanje c++ in pridobiti nova znanja na področju programiranja z kazalniki
    \item 
  \end{enumerate}
    \subsection{Uporabljena orodja}
    \begin{enumerate}
      \item nvim
      \item gcc in kasneje mingw
      \item spletni odpravljevalec napak (ang. debuger)%\cite{debuger} 
    \end{enumerate}
      \subsubsection{nvim}
      \subsubsection{gcc in mingw}
      \subsubsection{online gdb}
  \section{Programiranje}
    \subsection{Obrazložitev programa}
    Za lažjo delo sem program razdelil na tako imenovano ``glavno datoteko'', datoteko ``main'' in 3 dodatne datoteke z različnimi funkcijami. V nadaljevanju bom obrazložil delovanje programa najprej z vidika datoteke ``main'', kar predstavlja površinski pogled v delovanje programa, kasneje pa še vsako od funkcij iz vsake datoteke posebej.
      \subsubsection{config.h}
      datoteka ``config.h'' je ``glavna dototeka'', ki se uporablja za konsistenco definicij med različnimi datotekami\cite{cpplang}.\\ ``Glavno datoteko'' sem uporabil v naslednje namene:
      \begin{enumerate}
        \item definicija konstant, kot so: barve, ``razločevalci'' in daljša besedila,
        \item globalna definicija funkcij, ki jim v drugih datotekah priredim kodo, 
        \item globalna definicija struktur ``škatle'' in ``predmeta''.
      \end{enumerate}
      \subsubsection{main.cpp}
      Datoteka ``main'' je glavna datoteka, v kateri so definirane globalne spremenljivke, ki se uporabljajo čez celoten program. Dodatno v datoteki ``main'' živi tudi glavna, istoimenska funkcija, iz katere se kličejo vse nadaljne in, ki program zaključi. Poleg glavne funkcije pa se v tej datoteki nahaja še funkcija ``action\_switch'', ki na podlagi uporabniškega vnosa požene pravo funkcijo za nadaljevanje programa in izvedbo želenega ukaza.

      Program se iz vidika datoteke ``main'' izvaja po sledečem vrstnem redu:\\
      \begin{enumerate}
        \item dodeljen je spomin za globalne spremenljivke
        \item začne se funkcija ``main''
        \item program uporabnika vpraša za vnos ukaza, ki je le prva vpisana črka, iz seznama:
          \begin{itemize}
            \item d za ``delete'' ozeroma izbriši
            \item f za ``file'', ki uporabnika povpraša po imenu datoteke za odprtje
            \item h za ``help'' ozeroma pomoč, ki izpiše seznam ukazov in definicije
            \item l za ``list'' ozeroma naštej, ki, v primeru da je druga črka `a' izpiše atribute podanega predmeta ali škatle, v primeru da je druga črka `s' našteje vse predmete in škatle znotraj ali ob trenutni škatli, ter v primeru da druga črka ni nič od naštetega izpiše vse škatle in predmete v trenutni škatli
            \item m za ``move'' ozeroma premakni, ki spremeni trenutno škatlo v tisto podano po 3. mestu uporabniškega vhoda, kjer je ukaz ``mb'' rezerviran za premik nazaj v prejšnjo škatlo
            \item n za ``new'' ozeroma nov, ki v trenutni škatli, če je drugi znak vhoda `o' ustvari nov predmet, če pa ne pa novo škatlo
            \item s za ``search'', ki po trenutni škatli išče in izpiše predmete ali škatle z imenom ki se začne na niz podan po 2. znaku ukaza.
          \end{itemize}
        \item uporabnik izbere datoteko(drugega, razen izhoda, ne more, saj ni naložena še nobena datoteka)
        \item datoteka se prebere in v spomin se naložijo škatle in predmeti
        \item uporabnik lahko zdaj dodaja nove, jih birše, išče... 
        \item ko je uporabnik zadovoljen samo izvede ukaz `q', ki škatle in predmete shrani na isto datoteko, pri čemer izbriše prejšnje podatke, in program se zaključi
      \end{enumerate}
      \subsubsection{r\_w.cpp}
      Datoteka ``r\_w'' vsebuje funkcije za delo z datotekami in sicer:
      \begin{itemize}
        \item ``makefile'', ki prejme ime in preveri, da datoteka s tem imenom obstaja, če ne, jo naredi in vanjo vpiše osnovno strukturo datoteke
        \item ``writefile'', ki v predhodno odprto datoteko vpiše strukturo s pomočjo funkcije ``all\_box\_to\_file''
        \item ``readfile'', ki požene branje datoteke s pomočjo funkcije ``master\_switch''
        \item ``obj\_to\_file'' in ``box\_to\_file'', ki v datoteko vpišesta po en predmet in škatlo
        \item ``all\_box\_to\_file'', ki z uporabo prejšnjih dveh funkcij v datoteko zapiše vse povezane škatle in predmete
        \item ``print\_obj'' in ``print\_box'', ki izpišesta lastnosti predmeta in škatle
        \item ``print\_all\_box'', ki izpiše vse povezane škatle in predmete na standardni izhod
        \item ``get\_description'', ki iz vrstice z opisom predmeta ali škatle izlušči opis
        \item ``get\_sep\_spacing'', ki na podana naslova zapiše število znakov do dveh zaporednih pojavitev podanega znaka v vrstici
        \item ``get\_new\_line'', ki na naslov vrstice prepiše novo vrstico iz datoteke in izbriše staro
        \item ``insert\_corners'', ki na naslov ``ogljišč'' iz datoteke vpiše tri ``float-e'' (števila z vejico)
        \item ``master\_switch'', ki na podlagi prvega znaka v vrstici nadzira katere funkcije se bodo izvedle, s čimer omogoči, da se preberejo različne strukture
      \end{itemize}
      \subsubsection{ui.cpp}
      datoteka ``ui'' vsebuje funkcije povezane z uporabniškim umesnikom in ukazi, ki jih uporabnik lahko izvede. Funkcije so sledeče:
      \begin{itemize}
        \item ``delete\_box'', ki izbriše podano škatlo, njene predmete in vse njene podškatle
        \item ``delete\_with\_name'', ki na podlagi imena podanega poleg ukaza `d' izbriše škatlo, ozeroma, če ta ninajdena, predmet, s podanim imenom
        \item ``pritn\_path\_b'', ki rekurzivno izpiše pot do, s kazalnikom, podane škatle
        \item ``print\_path\_o'', ki izpiše pot do predmeta pri čemer uporabi prejšnjo funkcijo za zapis poti škatel
        \item ``print\_searched\_under'', ki na podlagi imen izpiše poti do vseh predmetov in škatel, katerih ime se začne z podanim nizom znakov
        \item ``print\_stats\_o'' in ``print\_stats\_b'', ki izpišeta atribute predmeta in škatle
        \item ``print\_atributes'', ki s pomočjo prejšnjih dveh funkcij izpiše atribute predmeta ali škatle z podanim imenom
        \item ``print\_names'', ki izpiše imena predmetov in škatel v trenutni škatli
        \item ``new\_object'' in ``new\_box'', ki uporabnika povprašata po podatkih in iz njih ustvarita nov predmet ali novo škatlo v trenunti
      \end{itemize}
      \subsubsection{linked\_funstions.cpp}

    \subsection{Potek programiranja}
      \subsubsection{Težave med programiranjem}
  \section{Možne izboljšave}
  \subsection{Kratkoročne}
  \subsection{Dolgoročne}
  \section{Zaključek}
  \newpage
  \section{Viri}
  \printbibliography[heading=none]

\end{document}



  %test text
  %\subsection{math}
  %test inline $F=ma$

  %test $\sum F = F_a + F_b$

  %\begin{equation}
  %W_p = \frac{1}{2} k s^2
  %\end{equation}
  %\newpage



  %\subsection{aj dont now}



  %\begin{figure}[H]
  %  \centering
  %  \includegraphics[width=10cm]{test_img.png}
  %  \caption{test}
  %  \label{fih:my_label}
  %\end{figure}
